\section{Phương pháp đề xuất}
\label{sec::phuong-phap-de-xuat}

\subsection{Nguyên lý}

Mục tiêu của bài báo muốn xây dựng quy trình tạo ra khuôn mặt được bảo vệ đáp ứng được đầy đủ chuẩn ISO/IEC24745 và tạo ra được khuôn mặt tự nhiên. Khác với các phương pháp che dấu danh tính khác, tác giả tận dụng đưa khuôn mặt vào trong không gian ẩn của các mô hình tạo sinh đã được tiền huấn luyện (cụ thể là StyleGAN), thực hiện biến đổi dựa trên tín hiệu đưa vào (còn gọi là chìa khóa). Để khuôn mặt được tự nhiên và sắc nét, người ta thiết kế ba hàm mất mát để điều chỉnh kết quả.

Không gian tiềm ẩn StyleGAN gồm 18 tầng xác định thông tin các mức ngữ nghĩa khác nhau. Tác giả mong muốn thay đổi tầng 3-7, nơi biểu diễn các đặc trưng gắn liền với danh tính, trong khi đó vẫn giữ nguyên các tầng 0-2 (kiểm soát cấu trúc tổng thể thô như tư thế đầu, hình dáng khuôn mặt) và tầng 8-17 (lưu giữ các chi tiết mịn như tóc, phông nền, ánh sáng). Mô hình biến đổi khóa có thể hủy, khởi tạo khóa thành mã biểu diễn thành một cấu trúc 18 tầng trong không gian của StyleGAN. Sau đó đưa thông tin đã khởi tạo của tầng 3-7 thay thế tầng chứa thông tin định danh của cấu trúc biểu diễn danh tính thật. Tác giả tận dụng kỹ thuật Style Mixing, mô hình tự động dung hòa thành một khuôn mặt hợp lý về mặt giải phẫu và ánh sáng.

Ảnh sinh ra đôi khi vẫn hơi sượng và mất biểu cảm, vì thế hệ thống tác giả đề xuất chạy thêm 50 vòng lặp để tối ưu. Bài báo sử dụng ba hàm mất mát, trong đó có hai hàm tăng khả năng nhận diện, hàm còn lại có chức năng kéo các chi tiết như nụ cười, hướng nhìn cho giống với ảnh gốc.

\subsection{Phát biểu bài toán}

Các không gian và thành phần chính của hệ thống định nghĩa như sau:
\begin{itemize}
    \item \textbf{Các không gian biểu diễn:} $\mathcal{X}$ là không gian ảnh khuôn mặt; $\mathcal{W}^+$ là không gian tiềm ẩn có thể kiểm soát của StyleGAN2; và $\mathcal{K}$ là không gian khóa.
    \item \textbf{Bộ mã hóa đảo ngược $\mathbf{E} : \mathcal{X} \rightarrow \mathcal{W}^+$:} sử dụng mạng tiền huấn luyện e4e (\textit{encoder for editing}) --- nhờ cân bằng tốt giữa độ trung thực tái tạo và khả năng chỉnh sửa trong $\mathcal{W}^+$ --- để ánh xạ ảnh khuôn mặt thật $\mathbf{x}_r \in \mathcal{X}$ sang mã tiềm ẩn $\mathbf{z}_r \in \mathcal{W}^+$. % TODO: bổ sung \cite{...} cho e4e (Tov et al., 2021) khi hoàn thiện references.bib
    \item \textbf{Bộ sinh ảnh $\mathbf{D} : \mathcal{W}^+ \rightarrow \mathcal{X}$:} mô hình StyleGAN2 đã tiền huấn luyện, nhận mã tiềm ẩn trung gian (ví dụ $\mathbf{z}_f \in \mathcal{W}^+$) để sinh ảnh khuôn mặt tổng hợp tương ứng ($\mathbf{x}_f \in \mathcal{X}$).
    \item \textbf{Hàm ánh xạ khóa $\mathbf{S} : \mathcal{K} \rightarrow \mathcal{W}^+$:} dùng khóa $\mathbf{k} \in \mathcal{K}$ làm hạt giống ngẫu nhiên (\textit{random seed}) để lấy mẫu mã tiềm ẩn $\mathbf{z}_f$.
    \item \textbf{Hàm hòa trộn $\mathbf{M} : \mathcal{W}^+ \times \mathcal{W}^+ \rightarrow \mathcal{W}^+$:} kết hợp mã tiềm ẩn ảnh thật $\mathbf{z}_r$ với mã tiềm ẩn từ khóa $\mathbf{z}_f$ thành mã tiềm ẩn trộn $\mathbf{z}_p$.
\end{itemize}

Từ các thành phần trên, ý tưởng cốt lõi của FaceAnonyMixer là tổng hợp ảnh khuôn mặt được bảo vệ $\mathbf{x}_p \in \mathcal{X}$ theo biểu thức:
\begin{equation}
\mathbf{x}_p = \mathbf{D}(\mathbf{z}_p), \quad \text{với } \mathbf{z}_p = \mathbf{M}(\mathbf{z}_r, \mathbf{z}_f), \quad \mathbf{z}_r = \mathbf{E}(\mathbf{x}_r), \quad \text{và } \mathbf{z}_f = \mathbf{S}(\mathbf{k}).
\end{equation}

Từ mã tiềm ẩn khởi tạo sau khi trộn $\mathbf{z}_p$, bài toán được quy về quá trình tìm kiếm mã tiềm ẩn tối ưu $\mathbf{z}_p^* \in \mathcal{W}^+$ (trong khi giữ cố định trọng số của $\mathbf{E}$ và $\mathbf{D}$) bằng cách cực tiểu hóa hàm mất mát đa mục tiêu:
\begin{equation}
\mathbf{z}_p^* = \operatorname*{arg\,min}_{\mathbf{z}_p \in \mathcal{W}^+} \mathcal{L}_{\text{total}} = \operatorname*{arg\,min}_{\mathbf{z}_p \in \mathcal{W}^+} \left( \lambda_1 \mathcal{L}_{\text{anon}} + \lambda_2 \mathcal{L}_{\text{idp}} + \lambda_3 \mathcal{L}_{\text{attr}} \right),
\end{equation}
trong đó $\lambda_1, \lambda_2, \lambda_3$ là các siêu tham số điều chỉnh mức độ đóng góp của từng thành phần mất mát:
\begin{itemize}
    \item $\mathcal{L}_{\text{anon}}$ (\textit{Anonymity Loss}): cực tiểu hóa độ tương đồng giữa đặc trưng danh tính của ảnh bảo vệ $\mathbf{x}_p$ và ảnh gốc $\mathbf{x}_r$, nhằm đảm bảo ẩn danh, không thể liên kết với danh tính thật và không thể đảo ngược.
    \item $\mathcal{L}_{\text{idp}}$ (\textit{Identity Preservation Loss}): cực tiểu hóa khoảng cách đặc trưng danh tính giữa các ảnh bảo vệ sinh ra từ nhiều biến thể khác nhau của cùng một người dưới cùng một khóa $\mathbf{k}$, nhằm bảo toàn tính nhất quán danh tính phục vụ xác thực.
    \item $\mathcal{L}_{\text{attr}}$ (\textit{Attribute Preservation Loss}): bảo toàn các thuộc tính phi danh tính (tư thế, biểu cảm, ánh sáng) giữa $\mathbf{x}_p$ và $\mathbf{x}_r$, nhằm duy trì độ tự nhiên thị giác và tính hữu dụng cho các tác vụ hạ nguồn.
\end{itemize}
Cuối cùng, ảnh khuôn mặt được bảo vệ hoàn chỉnh thu được thông qua phép giải mã $\mathbf{x}_p = \mathbf{D}(\mathbf{z}_p^*)$, qua đó đáp ứng trọn vẹn bốn yêu cầu của sinh trắc học có thể hủy theo chuẩn ISO/IEC 24745 mà không cần huấn luyện lại mô hình khi thay đổi khóa $\mathbf{k}$.
